%% NEW SECTION: A Different Computational Path to the Zeta Zeros
%% Draft for paper integration

\section{The Farey Spectroscope}\label{sec:spectroscope}

The explicit formula connects the per-step discrepancy to the nontrivial
zeros of~$\zeta(s)$.  In this section we exploit that connection in a
new direction: rather than \emph{assuming} the zeros and predicting
$\Delta W(p)$, we use $\Delta W(p)$ data to \emph{detect} the zeros.

\subsection{The Farey spectral function}

\begin{definition}\label{def:spectral-fn}
Given the insertion-deviation correlation ratio $R_2(p)$ for qualifying
primes $p$ (those with $M(p)\le-3$), define the \emph{Farey spectral function}
\[
  F(\gamma) \;=\; \Bigl|\,\sum_{p\;\text{qualifying}} R_2(p)\,p^{-1/2-i\gamma}\Bigr|^2.
\]
\end{definition}

The function $F(\gamma)$ acts as a periodogram of the signal $R_2(p)\,p^{-1/2}$
in the variable $\log p$.  Peaks in $F$ reveal frequencies present in the
oscillation of the Farey discrepancy across primes.

\begin{proposition}[GRH]\label{prop:spectroscope}
Assume the Riemann Hypothesis and that all nontrivial zeros of~$\zeta(s)$
are simple.  Then $F(\gamma)$ has local maxima at each~$\gamma_k$
(the imaginary parts of the nontrivial zeros), with peak heights satisfying
\[
  F(\gamma_k) \;\propto\; \left|\frac{1}{\rho_k\,\zeta'(\rho_k)}\right|^2
  \quad\text{as}\quad N\to\infty,
\]
where the proportionality constant depends on the number of primes used
and the sum $\sum_{p\le N}p^{-1}$.
\end{proposition}

\begin{proof}[Proof sketch]
Under GRH, the Farey counting error admits a spectral expansion
$E(N) \sim \sum_\rho N^\rho/(\rho\,\zeta'(\rho))$.  Since $R_2(p)$
captures the leading oscillatory component at prime orders, we have
$R_2(p)\,p^{-1/2}\approx\sum_k a_k\,p^{i\gamma_k}$ with
$a_k = 1/(\rho_k\,\zeta'(\rho_k))$.  Substituting into $F(\gamma)$:
at $\gamma = \gamma_j$ the $k{=}j$ term contributes constructively
($\sum_p p^{i(\gamma_j-\gamma_j)}=\sum_p 1 \sim N/\log N$),
while terms with $k\ne j$ involve oscillatory sums
$\sum_p p^{i(\gamma_k-\gamma_j)}$ that remain bounded under GRH.
\end{proof}

\subsection{Detection of $\gamma_1$}

Using $R_2(p)$ values for $N=2{,}729$ qualifying primes with $p\le56{,}209$,
we compute $F(\gamma)$ for $\gamma\in[5,52]$.

\begin{observation}\label{obs:gamma1}
$F(\gamma)$ attains its global maximum at $\gamma\approx 14.05$,
in agreement with the first nontrivial zero $\gamma_1=14.1347\ldots$
to within $0.6\%$.  A secondary peak appears at
$\gamma\approx 30.38$, matching $\gamma_4=30.4249\ldots$ to
$0.15\%$.
\end{observation}

Figure~\ref{fig:spectroscope} displays $F(\gamma)$ alongside the
known zero locations, and Figure~\ref{fig:juxtaposition} juxtaposes
the Farey spectral function with the classical Hardy~$Z$-function.

\subsection{Amplitude matching}

Proposition~\ref{prop:spectroscope} predicts that the \emph{relative}
peak heights $F(\gamma_k)/F(\gamma_1)$ should match
$|c_k/c_1|^2$ where $c_k = 1/(\rho_k\,\zeta'(\rho_k))$.

\begin{observation}\label{obs:amplitudes}
The correlation between predicted $|c_k/c_1|^2$ and observed
$F(\gamma_k)/F(\gamma_1)$ across the first ten zeros is
$r = 0.953$.  For $k=2{,}3$, the ratios agree to within $20{-}37\%$;
for $k\ge4$, resolution is limited by the number of primes available.
\end{observation}

\subsection{Extension to Dirichlet $L$-functions}

The spectroscope generalises to detect zeros of $L$-functions by
twisting with a Dirichlet character.

\begin{definition}\label{def:twisted}
For a Dirichlet character $\chi$ modulo~$q$, the \emph{$\chi$-twisted
Farey spectral function} is
\[
  F_\chi(\gamma) \;=\; \Bigl|\,\sum_{p} \chi(p)\,R_2(p)\,p^{-1/2-i\gamma}\Bigr|^2.
\]
\end{definition}

\begin{observation}\label{obs:dirichlet}
With $\chi_4$ (the non-principal character modulo~$4$),
$F_{\chi_4}(\gamma)$ shows peaks near
$\gamma'\approx 6.16$ and $\gamma'\approx 16.13$,
consistent with the first and fourth zeros of $L(s,\chi_4)$
(known values $\gamma'_1\approx 6.02$, $\gamma'_4\approx 16.00$).
Using the Kronecker symbol $\chi_3 = \bigl(\frac{\cdot}{3}\bigr)$
detects a zero of $L(s,\chi_3)$ near $\gamma\approx 16.40$
($\gamma''_3\approx 16.20$).
\end{observation}

Since the Dedekind zeta function of $\Q(i)$ factors as
$\zeta_{\Q(i)}(s) = \zeta(s)\cdot L(s,\chi_4)$,
the untwisted and $\chi_4$-twisted spectroscopes together
detect \emph{all} zeros of $\zeta_{\Q(i)}$.

\subsection{Convergence and comparison}

We compare the number of ``terms'' each method requires
to locate $\gamma_1$ to within $1\%$:

\begin{center}
\begin{tabular}{lcc}
\toprule
Method & Terms/primes & Error on $\gamma_1$ \\
\midrule
Farey spectroscope & 200 primes & $0.34\%$ \\
Riemann--Siegel (naive truncation) & 20 terms & $0.67\%$ \\
\bottomrule
\end{tabular}
\end{center}

The methods are comparable in raw convergence rate, but operate on
fundamentally different data: Riemann--Siegel evaluates $\zeta(s)$
directly, while the Farey spectroscope uses only Farey discrepancy
data from qualifying primes.  This independence makes the
spectroscope a potential verification tool.

\subsection{Pair correlation}

The autocorrelation of $F(\gamma)$ reveals structure at zero-pair
differences $\gamma_k - \gamma_l$, connecting to the Montgomery pair
correlation conjecture.

\begin{observation}\label{obs:pair}
The autocorrelation of $F(\gamma)$ peaks at lag $\approx 6.71$,
consistent with $\gamma_2 - \gamma_1 = 6.89$.
\end{observation}
